\section*{Capítulo \ref{ch:b3:environmental-toxicology} --- Química ambiental e toxicologia}

\begin{solution}{exo:b3:environmental-toxicology:1}
Uma diatômica homonuclear tem uma única vibração, que mantém nulo o momento dipolar,
e o argônio não tem vibração: nenhum deles pode absorver radiação infravermelha. O $\ce{CO2}$ absorve por
seu modo de deformação angular (e por seu estiramento antissimétrico), que criam um dipolo
oscilante; a deformação angular fica no meio da emissão da Terra.
\end{solution}

\begin{solution}{exo:b3:environmental-toxicology:2}
$10^{0.1} = 1.26$: $[\ce{H+}]$ aumenta 26\,\%.
\end{solution}

\begin{solution}{exo:b3:environmental-toxicology:3}
$(2.0 + 0.50)\ \unit{mmol/L} \times \qty{100.1}{mg/mmol} = \qty{250}{mg/L}$ de $\ce{CaCO3}$.
\end{solution}

\begin{solution}{exo:b3:environmental-toxicology:4}
N: $0.60/14.0 = \qty{0.043}{mmol/L}$; P: $0.010/31.0 = \qty{3.2e-4}{mmol/L}$; N\,:\,P $= 133$, muito acima de 16: o fósforo se esgota
primeiro e limita o crescimento.
\end{solution}

\begin{solution}{exo:b3:environmental-toxicology:5}
$L_0 = \mathrm{BOD_5}/(1 - \eu^{-0.23 \times 5}) = 170/0.683 = \qty{250}{mg/L}$.
\end{solution}

\begin{solution}{exo:b3:environmental-toxicology:6}
$\qty{4.40e-3}{L} \times \qty{0.0200}{mol/L} = \qty{8.80e-5}{mol}$ em \qty{0.1000}{L}: \omterm{def:b3:environmental-toxicology:alkalinity}{alcalinidade} $\qty{8.80e-4}{mol/L}$; como $\ce{HCO3-}$ (\qty{61.0}{g/mol}),
\qty{53.7}{mg/L}.
\end{solution}

\begin{solution}{exo:b3:environmental-toxicology:7}
$K_d = 0.020 \times 300 = \qty{6.0}{L/kg}$. Sorvido/total $= K_dm_s/(K_dm_s + V_w) = 9.0/(9.0 + 0.30) = 0.97$: 97\,\% sorvido.
\end{solution}

\begin{solution}{exo:b3:environmental-toxicology:8}
Benzeno: $\log\mathrm{BCF} = 0.85 \times 2.13 - 0.70 = 1.11$, BCF de cerca de 13. PCB: $0.85 \times 6.67 - 0.70 = 4.97$, BCF de cerca de
\num{9e4}: ele se concentra nos peixes cerca de sete mil vezes mais que o benzeno; com um
$K_{ow}$ desses, a regressão está perto de seu limite de validade.
\end{solution}

\begin{solution}{exo:b3:environmental-toxicology:9}
$\qty{192}{mg/kg} \times \qty{70}{kg} = \qty{13}{g}$, cerca de 150 xícaras. As espécies diferem na absorção e no metabolismo,
e uma $\mathrm{LD_{50}}$ não é um limiar seguro; os efeitos nocivos começam muito abaixo dela.
\end{solution}

\begin{solution}{exo:b3:environmental-toxicology:10}
PNEC $= \qty{0.50}{mg/L}/1000 = \qty{0.50}{\micro g/L}$; $\mathrm{RQ} = 0.20/0.50 = 0.40$, abaixo de 1: nenhuma preocupação com essa exposição.
\end{solution}

\begin{solution}{exo:b3:environmental-toxicology:11}
Plâncton/água: $\qty{0.015}{mg/kg}/\qty{2.0e-6}{mg/L} = \qty{7.5e3}{L/kg}$ (bioconcentração). Peixe/plâncton: $0.30/0.015 = 20$;
ave/peixe: $6.0/0.30 = 20$ (\omterm{def:b3:environmental-toxicology:bio}{biomagnificação} em cada etapa); ave/água: $\num{3e6}$.
\end{solution}

\begin{solution}{exo:b3:environmental-toxicology:12}
$\qty{1.0e9}{g}/\qty{64.1}{g/mol} = \qty{1.56e7}{mol}$ de $\ce{SO2}$, que dão outros tantos mols de $\ce{H2SO4}$: $\qty{1.53e9}{g}$, cerca de \qty{1500}{t}. Isso
libera $\qty{3.12e7}{mol}$ de $\ce{H+}$; a pH 4.0, $[\ce{H+}] = \qty{1.0e-4}{mol/L}$: $\qty{3.1e11}{L}$ de chuva.
\end{solution}

\begin{solution}{pb:b3:environmental-toxicology:1}
\textbf{1.} Ele é dez mil vezes mais solúvel no octanol que na água: hidrofóbico,
vai se sorver na matéria orgânica e se acumular na gordura.

\textbf{2.} $K_{oc} = 0.41 \times 10^4 = \qty{4100}{L/kg}$; $K_{sw} = f_{oc}K_{oc} = 0.040 \times 4100 = \qty{164}{L/kg}$.

\textbf{3.} $K_{fw} = 0.048 \times 10^4 = \qty{480}{L/kg}$.

\textbf{4.} Não: no equilíbrio, o ar contém $10^{-4}$ da concentração na água.

\textbf{5.} A molécula neutra e hidrofóbica se dissolve na matéria orgânica, como no
octanol; as superfícies minerais são polares e recobertas de água.

\textbf{6.} O sedimento, cuja capacidade, $K_{sw}m_s$, é a maior.

\textbf{7.} $V_w = \qty{1.0e9}{L}$; $K_{aw}V_a = \qty{1.0e8}{L}$; $K_{sw}m_s = \qty{1.64e10}{L}$; $K_{fw}m_f = \qty{4.8e6}{L}$.

\textbf{8.} $C_w = \qty{1.0e8}{mg}/\qty{1.75e10}{L} = \qty{5.7e-3}{mg/L}$, \qty{5.7}{\micro g/L}.

\textbf{9.} Água \qty{5.7}{kg}, ar \qty{0.57}{kg}, sedimento \qty{94}{kg}, peixes \qty{0.027}{kg}.

\textbf{10.} Ar $\qty{5.7e-7}{mg/L}$; sedimento \qty{0.94}{mg/kg}; peixes \qty{2.7}{mg/kg}.

\textbf{11.} O equilíbrio entre todos os compartimentos ao mesmo tempo, sem degradação, sem
entrada nem saída, e com compartimentos bem misturados.

\textbf{12.} Com degradação e fluxos em estado estacionário (nível II), o estoque é fixado pelas
velocidades de entrada e de perda; com velocidades de transferência entre compartimentos (nível III),
as concentrações deixam de estar em razões de equilíbrio, e o compartimento que
recebe a entrada fica enriquecido.

\textbf{13.} $k = \ln 2/\qty{60}{d} = \qty{0.0116}{d^{-1}}$.

\textbf{14.} $\eu^{-0.0116 \times 365} = 0.0145$: restam 1.5\,\%.

\textbf{15.} $\qty{5.7}{\micro g/L} \times 0.0145 = \qty{0.083}{\micro g/L}$.

\textbf{16.} A degradação é muitas vezes mais lenta no sedimento frio e anóxico, e o pesticida
sorvido é liberado de volta para a água lentamente, mantendo-o presente depois que a
água já foi depurada.

\textbf{17.} Ligações que resistem à hidrólise, à oxidação e ao ataque microbiano (C--Cl
aromáticas, por exemplo), baixa solubilidade em água e uma sorção que o protege dos
degradadores.

\textbf{18.} PNEC $= \qty{50}{\micro g/L}/10 = \qty{5.0}{\micro g/L}$.

\textbf{19.} $\mathrm{RQ} = 5.7/5.0 = 1.1$.

\textbf{20.} Pouco acima de 1: indica-se um risco, e a avaliação precisa ser refinada (melhores
dados de toxicidade, concentrações medidas) ou a exposição reduzida.

\textbf{21.} \qty{2.7}{mg/kg}, quase três vezes o limite de \qty{1.0}{mg/kg}: os peixes não devem ser consumidos
até que a concentração caia.

\textbf{22.} As concentrações na água, no sedimento e nos peixes ao longo do tempo e em vários pontos,
para verificar o modelo e acompanhar a diminuição.

\textbf{23.} Menos pesticida aplicado, aplicação longe das chuvas e da margem,
faixas de proteção que retêm o escoamento, ou uma alternativa menos persistente.

\textbf{24.} Ele divide o nível sem efeito medido para cobrir a distância entre algumas poucas
espécies de laboratório e um ecossistema inteiro; um fator de 10 em vez de 1000 muda a
PNEC, e o veredito, por um fator de cem.

\textbf{25.} O \omterm{def:b3:environmental-toxicology:ecotoxicology}{quociente de risco} PEC/PNEC é de cerca de 1.1, pouco acima de 1.
\end{solution}
